\kchapter{Kubernetes in Production}{control-plane}
\label{ch:production}

\begin{objectives}
\begin{itemize}
  \item managed vs.\ self-managed clusters;
  \item high availability for the control plane and your apps;
  \item upgrades without drama;
  \item backup and disaster recovery with etcd snapshots and Velero;
  \item multi-cluster, service mesh, cost management, and a go-live checklist.
\end{itemize}
\end{objectives}

\section{Managed or self-managed?}

\begin{center}
\ktablestart
\begin{tabularx}{\linewidth}{P{3.0cm} L L}
\kheader \kth{} & \kth{Managed (EKS, GKE, AKS\ldots)} & \kth{Self-managed (kubeadm, k3s, RKE2, Talos)} \\
Control plane & Run, patched, and backed up by the provider. & You run etcd, API servers, certificates. \\
Upgrades & A button or API call (nodes still need care). & Fully your responsibility. \\
Cost & Small control-plane fee. & Engineers' time. \\
When & Almost always, in the cloud. & On-premises, edge, air-gapped, special compliance. \\
\end{tabularx}
\end{center}

\begin{keyidea}
Unless you have a strong reason, \textbf{use a managed Kubernetes service}.
Your competitive advantage is your application, not running etcd.
\end{keyidea}

\section{High availability}

\subsection{Control plane}

\begin{itemize}
  \item 3 (or 5) control-plane nodes across availability zones;
  \item etcd with an odd number of members (quorum = majority);
  \item API servers behind a load balancer.
\end{itemize}

\subsection{Your applications}

\begin{itemize}
  \item at least \textbf{2--3 replicas} for anything user-facing;
  \item \textbf{topology spread} across zones and nodes;
  \item \textbf{PodDisruptionBudgets} so drains never take everything down;
  \item \textbf{readiness probes} and graceful shutdown for zero-downtime deploys;
  \item \textbf{requests/limits} so noisy neighbours cannot starve you;
  \item node pools across multiple zones.
\end{itemize}

\section{Upgrades}

Kubernetes releases a minor version about every four months and supports each
for roughly 14 months. Falling behind is painful; upgrade regularly.

\begin{enumerate}
  \item \textbf{Read the release notes} and the \emph{deprecated API} guide.
        Check your manifests and charts with tools like \textbf{pluto} or
        \textbf{kubent}.
  \item Upgrade \textbf{one minor version at a time} (1.32 $\rightarrow$ 1.33 $\rightarrow$ 1.34).
  \item Upgrade the \textbf{control plane first}, then nodes. Nodes may lag
        the control plane by a few minor versions, never lead it.
  \item Upgrade nodes by \textbf{cordon \& drain} or by creating a new node pool
        and moving workloads (blue/green nodes).
  \item Upgrade add-ons (CNI, CSI, CoreDNS, Gateway controller) to compatible versions.
  \item Practise on staging first.
\end{enumerate}

\section{Backup and disaster recovery}

\subsection{etcd snapshots (self-managed clusters)}

\begin{bashcode}{terminal}
ETCDCTL_API=3 etcdctl snapshot save /backup/etcd-$(date +%F).db \
  --endpoints=https://127.0.0.1:2379 \
  --cacert=/etc/kubernetes/pki/etcd/ca.crt \
  --cert=/etc/kubernetes/pki/etcd/server.crt \
  --key=/etc/kubernetes/pki/etcd/server.key

etcdutl snapshot status /backup/etcd-2026-10-10.db --write-out=table
\end{bashcode}

\subsection{Velero: backing up workloads and volumes}

\textbf{Velero} backs up Kubernetes objects and persistent volume data to object
storage, and restores them --- into the same or a different cluster.

\begin{bashcode}{terminal}
velero backup create shop-daily --include-namespaces shop --snapshot-volumes
velero schedule create shop-nightly --schedule "0 3 * * *" --include-namespaces shop --ttl 720h
velero restore create --from-backup shop-daily
\end{bashcode}

\begin{warningbox}
A backup you have never restored is a hope, not a backup. Schedule regular
restore drills into a scratch cluster and measure how long they take (your
real Recovery Time Objective).
\end{warningbox}

Remember too that if \emph{everything} (manifests, Helm values, Argo CD
Applications) is in Git, rebuilding a cluster's configuration is ``create cluster,
install Argo CD, point at the repo''. Only \emph{data} needs separate backups.

\section{Multi-cluster}

Teams run several clusters for isolation (prod vs.\ non-prod), regions, latency,
or blast-radius reduction. Common tools:

\begin{itemize}
  \item \textbf{Argo CD ApplicationSets} / \textbf{Flux} to deploy to many clusters from one repo;
  \item \textbf{Cluster API} to create and upgrade clusters declaratively;
  \item global load balancing / DNS to route users to the nearest healthy cluster.
\end{itemize}

\section{Service mesh}

A \kicon{istio}~\kterm{service mesh} (Istio, Linkerd, Cilium Service Mesh) adds
a transparent network layer between services, giving you:

\begin{itemize}
  \item \textbf{mutual TLS} between every service, with identities, automatically;
  \item fine-grained \textbf{traffic control}: retries, timeouts, circuit breaking, canaries;
  \item uniform \textbf{telemetry}: golden metrics and traces for every call.
\end{itemize}

Modern meshes offer \emph{sidecar-less} modes (Istio Ambient, Cilium) that lower
the overhead. Adopt a mesh when you actually need mTLS everywhere or advanced
traffic control --- it is a significant piece of infrastructure.

\section{Cost management}

\begin{itemize}
  \item \textbf{Right-size requests} using real usage (VPA recommendations,
        Prometheus data). Over-requesting is the number-one source of waste.
  \item Use \textbf{autoscaling} at every level, and scale dev environments to zero at night.
  \item Run interruption-tolerant workloads on \textbf{spot/preemptible} nodes.
  \item Use \textbf{ARM nodes} where your images support them.
  \item Allocate costs per namespace/team with \textbf{OpenCost} or \textbf{Kubecost}.
\end{itemize}

\section{Go-live checklist}

\begin{multicols}{2}
\begin{itemize}
  \item[\textcolor{k8sblue}{\faCheckSquare}] Managed control plane, multi-zone nodes
  \item[\textcolor{k8sblue}{\faCheckSquare}] Everything in Git; GitOps sync
  \item[\textcolor{k8sblue}{\faCheckSquare}] Images pinned, scanned, signed
  \item[\textcolor{k8sblue}{\faCheckSquare}] Probes, requests/limits, PDBs
  \item[\textcolor{k8sblue}{\faCheckSquare}] HPA + cluster autoscaling
  \item[\textcolor{k8sblue}{\faCheckSquare}] Topology spread across zones
  \item[\textcolor{k8sblue}{\faCheckSquare}] TLS via cert-manager
  \item[\textcolor{k8sblue}{\faCheckSquare}] NetworkPolicies, restricted PSS
  \item[\textcolor{k8sblue}{\faCheckSquare}] Least-privilege RBAC, SSO
  \item[\textcolor{k8sblue}{\faCheckSquare}] External secret manager
  \item[\textcolor{k8sblue}{\faCheckSquare}] Metrics, logs, traces, SLO alerts
  \item[\textcolor{k8sblue}{\faCheckSquare}] Runbooks for every alert
  \item[\textcolor{k8sblue}{\faCheckSquare}] Backups + tested restores
  \item[\textcolor{k8sblue}{\faCheckSquare}] Upgrade plan and cadence
\end{itemize}
\end{multicols}

\begin{recap}
\begin{itemize}
  \item Prefer managed clusters; spend your time on your apps.
  \item HA = replicas + spread + PDBs + probes, across zones.
  \item Upgrade one minor version at a time, control plane first.
  \item Back up data (Velero, snapshots) and keep config in Git; test restores.
\end{itemize}
\end{recap}
